\chapter{Lars Valerian Ahlfors --- Fields Medal 1936}\label{ch1936larsvalerianahlfors:lars-valerian-ahlfors-fields-medal-1936-a-comprehensive-report}

\begin{center}\rule{0.5\linewidth}{0.5pt}\end{center}

\textbf{Prepared:} July 2026 \textbf{Subject:} Lars Valerian Ahlfors (b.
April 18, 1907, Helsinki, Finland) \textbf{Position:} Professor, Harvard
University \textbf{Doctoral Advisor:} Ernst Lindelöf and Rolf Nevanlinna
(University of Helsinki) \textbf{Fields Medal Citation:} \emph{``Awarded
medal for research on covering surfaces related to Riemann surfaces of
inverse functions of entire and meromorphic functions. Opened up new
fields of analysis.''}

\begin{center}\rule{0.5\linewidth}{0.5pt}\end{center}

\section{1. Introduction}\label{ch1936larsvalerianahlfors:introduction}

Lars Valerian Ahlfors was one of the two inaugural recipients of the
Fields Medal in 1936, honored for his substantial geometric approach
to complex analysis. His work transformed Nevanlinna's value
distribution theory from an analytically intricate subject into a
geometrically intuitive framework based on covering surfaces and
conformal metrics. This geometric reformation not only simplified the
proofs of Nevanlinna's main theorems but significantly extended their
reach beyond meromorphic functions to the broader class of
quasiconformal mappings.

\begin{prereqs} Complex numbers, analytic functions, contour integration, basic knowledge of Riemann surfaces. \end{prereqs}

\begin{objectives} Explain the Ahlfors five-disk theorem and its significance; describe the connection between curvature and covering spaces; state the main results of Ahlfors's geometric function theory. \end{objectives}


Ahlfors's contributions span the full spectrum of complex analysis:

\begin{itemize}
\tightlist
\item
 The \textbf{Ahlfors theory of covering surfaces}, providing a
 geometric reformulation of Nevanlinna theory through the language of
 Riemann surfaces and conformal metrics.
\item
 The \textbf{Denjoy--Carleman--Ahlfors theorem}, establishing the sharp
 upper bound for the number of asymptotic values of an entire function
 of finite order.
\item
 The \textbf{Ahlfors--Schwarz lemma}, generalizing Schwarz's classical
 lemma to Riemann surfaces with curvature bounds.
\item
 The \textbf{Ahlfors finiteness theorem}, proving that a finitely
 generated Kleinian group has a finite-sided fundamental polyhedron.
\item
 The \textbf{Ahlfors--Bers theorem} on the measurable Riemann mapping
 theorem, foundational to Teichmüller theory.
\item
 The theory of \textbf{extremal length} and \textbf{conformal
 invariants}, providing powerful tools for geometric function theory.
\item
 His influential textbook \emph{Complex Analysis}, which remains a
 classic in the field after over seven decades.
\end{itemize}

This report provides a comprehensive overview of Ahlfors's life, his
monumental contributions to mathematics, and the lasting impact of his
work.

\begin{center}\rule{0.5\linewidth}{0.5pt}\end{center}

\section{2. Biographical Background}\label{ch1936larsvalerianahlfors:biographical-background}

\subsection{2.1 Early Life and
Education}\label{ch1936larsvalerianahlfors:early-life-and-education}

Lars Valerian Ahlfors was born in Helsinki, Finland, on April 18, 1907.
His father was a professor of engineering at the Helsinki University of
Technology. Ahlfors showed early mathematical talent and entered the
University of Helsinki at the age of 17. There he came under the
influence of Ernst Lindelöf and Rolf Nevanlinna, two of Finland's most
distinguished mathematicians, who introduced him to the frontiers of
complex analysis.

At 21, Ahlfors followed Nevanlinna to the Eidgenössische Technische
Hochschule (ETH) in Zürich, where he began working at a research level.
There he produced his first major work on the asymptotic values of
entire functions, based on a new approach to conformal mapping.
Self-effacingly, Ahlfors credited Nevanlinna and George Pólya for their
considerable help; they in turn insisted he publish the results solely
in his own name.

\subsection{2.2 Academic Career}\label{ch1936larsvalerianahlfors:academic-career}

\begin{itemize}
\tightlist
\item
 \textbf{1928:} Master's degree, University of Helsinki.
\item
 \textbf{1930:} PhD, University of Helsinki, under Ernst Lindelöf and
 Rolf Nevanlinna.
\item
 \textbf{1933--1935:} Docent, University of Helsinki.
\item
 \textbf{1935--1938:} Visiting Lecturer, Harvard University.
\item
 \textbf{1936:} Fields Medal awarded at the ICM in Oslo.
\item
 \textbf{1938--1944:} Professor, University of Helsinki.
\item
 \textbf{1944--1946:} Professor, University of Zurich.
\item
 \textbf{1946--1977:} Professor, Harvard University.
\item
 \textbf{1977:} Professor Emeritus, Harvard University.
\item
 \textbf{1981:} Wolf Prize in Mathematics.
\item
 \textbf{1996:} Died in Helsinki at age 89.
\end{itemize}

\subsection{2.3 Doctoral Work}\label{ch1936larsvalerianahlfors:doctoral-work}

Ahlfors's doctoral dissertation, completed in 1930 under the supervision
of Lindelöf and Nevanlinna, dealt with asymptotic values of meromorphic
functions. This work already exhibited his characteristic geometric
approach and set the stage for his later development of covering surface
theory. The dissertation contained what later became known as the
Ahlfors--Carleman theorem on the number of asymptotic values of an
entire function of finite order.

\begin{center}\rule{0.5\linewidth}{0.5pt}\end{center}

\section{3. The Denjoy--Carleman--Ahlfors
Theorem}\label{ch1936larsvalerianahlfors:the-denjoycarlemanahlfors-theorem}

\subsection{3.1 Background}\label{ch1936larsvalerianahlfors:background}

An \textbf{asymptotic value} of an entire function \(f(z)\) is a value
\(a \in \mathbb{C}\) such that there exists a path \(\gamma(t)\) tending
to infinity with \(\lim_{t\to\infty} f(\gamma(t)) = a\). A classical
theorem of Émile Picard states that a non-constant entire function can
omit at most one finite value. The question of how many asymptotic
values an entire function can have was raised by Arnaud Denjoy in 1907.

\subsection{3.2 The Theorem}\label{ch1936larsvalerianahlfors:the-theorem}

Denjoy conjectured that an entire function of finite order \(\rho\) can
have at most \(2\rho\) asymptotic values (finite or infinite). Carleman
proved a special case. Ahlfors, in his 1930 paper, proved the complete
theorem:

\begin{quote}
\textbf{Denjoy--Carleman--Ahlfors Theorem:} An entire function of finite
order \(\rho\) has at most \(2\rho\) distinct asymptotic values.
\end{quote}

The proof uses the method of \textbf{extremal length} and conformal
mapping, showing that each asymptotic value corresponds to a path whose
pre-image in a certain covering surface has specific growth properties.

\subsection{3.3 Impact}\label{ch1936larsvalerianahlfors:impact}

The theorem was the first major result establishing a precise
relationship between the growth of an entire function and its asymptotic
behavior. It demonstrated the power of geometric methods in complex
analysis and established Ahlfors's reputation as a rising star.

\begin{center}\rule{0.5\linewidth}{0.5pt}\end{center}

\section{4. Nevanlinna Theory}\label{ch1936larsvalerianahlfors:nevanlinna-theory}

\subsection{4.1 Classical Nevanlinna
Theory}\label{ch1936larsvalerianahlfors:classical-nevanlinna-theory}

Rolf Nevanlinna's value distribution theory, developed in 1925, studies
the distribution of values of a meromorphic function
\(f: \mathbb{C} \to \overline{\mathbb{C}}\). The fundamental quantities
are:

\begin{itemize}
\tightlist
\item
 The \textbf{proximity function} \(m(r, f)\), which measures how close
 \(f(z)\) is to \(\infty\) on \(|z| = r\).
\item
 The \textbf{counting function} \(N(r, a)\), which counts how many
 times \(f(z) = a\) in \(|z| < r\).
\item
 The \textbf{characteristic function} \(T(r, f)\), which measures the
 growth of \(f\).
\end{itemize}

Nevanlinna's two main theorems state:

\begin{quote}
\textbf{First Main Theorem:} \(T(r, f) = m(r, a) + N(r, a) + O(1)\) for
any \(a \in \overline{\mathbb{C}}\).
\end{quote}

\begin{quote}
\textbf{Second Main Theorem:} For \(q \geq 3\) distinct values
\(a_1, \ldots, a_q\),
\[(q-2)T(r, f) \leq \sum_{j=1}^q N_1(r, a_j) + S(r, f)\] where
\(N_1(r, a)\) is the counting function with multiplicities reduced and
\(S(r, f)\) is a small error term.
\end{quote}

\subsection{4.2 Ahlfors's Geometric
Reformulation}\label{ch1936larsvalerianahlfors:ahlforss-geometric-reformulation}

Ahlfors recognized that Nevanlinna theory could be understood
geometrically. The key insight: consider
\(f: \mathbb{C} \to \overline{\mathbb{C}}\) as a \textbf{covering
surface} over the Riemann sphere. The characteristic function
\(T(r, f)\) measures the area of the image surface (the part of \(f\)'s
graph lying over the disc of radius \(r\)), while the counting function
measures how many sheets cover each point.

In this geometric framework, Nevanlinna's deficiency
\(\delta(a, f) = \liminf_{r\to\infty} \frac{m(r, a)}{T(r, f)}\) is
interpreted as the proportion of the sphere that is ``thinly covered''
near the value \(a\). The second main theorem becomes an isoperimetric
inequality on the covering surface, relating the area, boundary length,
and Euler characteristic.

Ahlfors's insight was that the central estimates of value distribution
theory --- which Nevanlinna had derived through intricate manipulations
of the Poisson--Jensen formula --- could be obtained instead from a
single application of the Gauss--Bonnet theorem combined with the
isoperimetric inequality. This made the proofs not only simpler but also
conceptually clearer.

\subsection{4.3 The 14-Page Paper}\label{ch1936larsvalerianahlfors:the-14-page-paper}

In 1935, Ahlfors published his geometric version of Nevanlinna theory in
just 14 pages. Carathéodory remarked at the Fields Medal ceremony that
it was difficult to say which was more surprising: that Nevanlinna could
develop his entire theory without the geometric picture, or that Ahlfors
could condense the whole theory into fourteen pages.

\begin{center}\rule{0.5\linewidth}{0.5pt}\end{center}

\section{5. The Theory of Covering
Surfaces}\label{ch1936larsvalerianahlfors:the-theory-of-covering-surfaces}

\subsection{5.1 Basic Concepts}\label{ch1936larsvalerianahlfors:basic-concepts}

A \textbf{covering surface} over a base surface \(S\) is a surface
\(\widetilde{S}\) together with a continuous surjection
\(p: \widetilde{S} \to S\) that is a local homeomorphism. In Ahlfors's
theory, the map \(f: \mathbb{C} \to \overline{\mathbb{C}}\) is viewed as
a covering surface of the Riemann sphere, possibly with branch points at
critical points of \(f\).

The central object is the \textbf{Ahlfors covering surface}
\(\Omega_f(r)\), defined as the part of the \(f\)-image of \(|z| < r\)
lifted to the appropriate branched covering. Its boundary consists of
arcs that project to curves on the Riemann sphere, and its interior
geometry encodes the value distribution of \(f\).

\subsection{5.2 Ahlfors's Main Idea}\label{ch1936larsvalerianahlfors:ahlforss-main-idea}

The guiding intuition: if a meromorphic function is given, the
fundamental quantities studied in Nevanlinna theory --- such as the
number of \(a\)-points in a disc of radius \(r\) --- can be interpreted
as geometric properties of the covering surface. Specifically:

\begin{itemize}
\tightlist
\item
 The \textbf{mean covering number} measures how many times, on average,
 the covering surface lies over each point of the base.
\item
 The \textbf{Euler characteristic} of the covering surface relates to
 the number of branch points.
\item
 The \textbf{isoperimetric inequality} on the covering surface gives
 the main estimates.
\end{itemize}

\subsection{5.3 The Two Main Theorems}\label{ch1936larsvalerianahlfors:the-two-main-theorems}

Ahlfors proved:

\begin{quote}
\textbf{Ahlfors's First Main Theorem:} For a covering surface \(\Omega\)
over the Riemann sphere, the mean covering number equals the Euler
characteristic of \(\Omega\) divided by \(2\pi\), up to boundary terms.
\end{quote}

\begin{quote}
\textbf{Ahlfors's Second Main Theorem:} If \(\Omega\) is a covering
surface with \(q\) boundary components that project to distinct points,
then \[(q-2)A \leq 2\pi\chi_{-} + B\] where \(A\) is the area,
\(\chi_{-}\) is the negative part of the Euler characteristic, and \(B\)
is a boundary term.
\end{quote}

\subsection{5.4 Significance}\label{ch1936larsvalerianahlfors:significance}

Ahlfors's covering surface theory was described by Carathéodory as ``a
much more radical departure from Nevanlinna's own methods'' than his
earlier geometric reformulation. It extended Nevanlinna theory from
meromorphic functions to a much wider class of maps:
\textbf{quasiconformal mappings} and, indeed, any map that can be
interpreted as a covering surface.

\begin{center}\rule{0.5\linewidth}{0.5pt}\end{center}

\section{6. Quasiconformal Mappings}\label{ch1936larsvalerianahlfors:quasiconformal-mappings}

\subsection{6.1 Definition}\label{ch1936larsvalerianahlfors:definition}

A \textbf{quasiconformal mapping} \(f: D \to D'\) between domains in
\(\mathbb{C}\) is a homeomorphism that preserves orientation and has
bounded dilatation. The \textbf{dilatation} \(K(f)\) measures the
distortion of angles:

\[K(f) = \sup_{z \in D} \frac{|f_z(z)| + |f_{\bar{z}}(z)|}{|f_z(z)| - |f_{\bar{z}}(z)|}\]

When \(K(f) = 1\), \(f\) is conformal. Quasiconformal mappings were
introduced by Grötzsch and developed by Teichmüller, but Ahlfors made
them into a central tool of complex analysis.

\subsection{6.2 Ahlfors's Contributions}\label{ch1936larsvalerianahlfors:ahlforss-contributions}

Ahlfors recognized that his covering surface theory applied naturally to
quasiconformal mappings. He showed:

\begin{itemize}
\tightlist
\item
 The Ahlfors--Schwarz lemma provides curvature bounds for
 quasiconformal metrics.
\item
 Quasiconformal mappings preserve the class of null-sets for extremal
 length.
\item
 The measurable Riemann mapping theorem (with Bers) shows that any
 Beltrami equation has a solution.
\end{itemize}

\subsection{6.3 The Measurable Riemann Mapping
Theorem}\label{ch1936larsvalerianahlfors:the-measurable-riemann-mapping-theorem}

The \textbf{Beltrami equation} is:

\[\frac{\partial f}{\partial \bar{z}} = \mu(z) \frac{\partial f}{\partial z}\]

where \(\mu\) is a measurable function with \(\|\mu\|_\infty < 1\). The
\textbf{Ahlfors--Bers theorem} (proved jointly with Lipman Bers) states:

\begin{quote}
\textbf{Measurable Riemann Mapping Theorem:} For any measurable \(\mu\)
with \(\|\mu\|_\infty < 1\), there exists a quasiconformal homeomorphism
\(f\) of the plane satisfying the Beltrami equation. Moreover, \(f\)
depends complex-analytically on the parameter \(\mu\).
\end{quote}

\subsection{6.4 Impact}\label{ch1936larsvalerianahlfors:impact-1}

The Ahlfors--Bers theorem is foundational for Teichmüller theory,
complex dynamics, and the theory of Kleinian groups. It shows that the
space of all quasiconformal mappings is a complex Banach manifold,
enabling the use of analytic methods in deformation theory.

\begin{center}\rule{0.5\linewidth}{0.5pt}\end{center}

\section{7. The Ahlfors--Schwarz Lemma}\label{ch1936larsvalerianahlfors:the-ahlforsschwarz-lemma}

\subsection{7.1 Classical Schwarz Lemma}\label{ch1936larsvalerianahlfors:classical-schwarz-lemma}

The classical Schwarz lemma states that a holomorphic map
\(f: \mathbb{D} \to \mathbb{D}\) with \(f(0) = 0\) satisfies
\(|f(z)| \leq |z|\) for all \(z \in \mathbb{D}\) and \(|f'(0)| \leq 1\).

\subsection{7.2 The Ahlfors
Generalization}\label{ch1936larsvalerianahlfors:the-ahlfors-generalization}

Ahlfors generalized this to Riemann surfaces with curvature bounds. The
\textbf{Ahlfors--Schwarz lemma} states:

\begin{quote}
\textbf{Ahlfors--Schwarz Lemma:} Let \(M\) be a Riemann surface equipped
with a conformal metric \(ds_M^2\) of curvature \(\leq -1\), and let
\(N\) be a Riemann surface with a conformal metric \(ds_N^2\) of
curvature \(\geq -1\). Then any holomorphic map \(f: M \to N\) is
distance-decreasing: \(f^*ds_N \leq ds_M\).
\end{quote}

\subsection{7.3 Proof via Curvature}\label{ch1936larsvalerianahlfors:proof-via-curvature}

The proof uses the \textbf{maximum principle} for the Laplacian of the
ratio \(f^*ds_N/ds_M\). If the ratio exceeds 1 at some point, then at
its maximum the curvature comparison gives a contradiction.

\subsection{7.4 Applications}\label{ch1936larsvalerianahlfors:applications}

The Ahlfors--Schwarz lemma has numerous applications:

\begin{itemize}
\tightlist
\item
 It provides a unified proof of Picard's theorem and other value
 distribution results.
\item
 It gives bounds on the growth of holomorphic maps between hyperbolic
 Riemann surfaces.
\item
 It is used in the study of holomorphic foliations and complex
 dynamics.
\item
 It generalizes to higher dimensions via the Ahlfors--Yau--Royden
 Schwarz lemma.
\end{itemize}

\begin{center}\rule{0.5\linewidth}{0.5pt}\end{center}

\section{8. Extremal Length}\label{ch1936larsvalerianahlfors:extremal-length}

\subsection{8.1 Definition}\label{ch1936larsvalerianahlfors:definition-1}

\textbf{Extremal length} is a conformal invariant introduced by Ahlfors
and Beurling in the 1940s. For a family \(\Gamma\) of curves in a domain
\(D\), the extremal length is:

\[\lambda(\Gamma) = \sup_{\rho} \frac{\inf_{\gamma \in \Gamma} \int_\gamma \rho \, ds}{\int_D \rho^2 \, dx\,dy}\]

where the supremum is taken over all non-negative Borel functions
\(\rho\) on \(D\).

\subsection{8.2 Properties}\label{ch1936larsvalerianahlfors:properties}

Extremal length is: - \textbf{Conformally invariant:}
\(\lambda(\Gamma) = \lambda(f(\Gamma))\) for any conformal map \(f\). -
\textbf{Monotone:} If \(\Gamma_1 \subset \Gamma_2\), then
\(\lambda(\Gamma_1) \geq \lambda(\Gamma_2)\). - \textbf{Countably
subadditive.} - \textbf{A metric on the space of curve families.}

\subsection{8.3 Ahlfors's Contributions}\label{ch1936larsvalerianahlfors:ahlforss-contributions-1}

Ahlfors and Beurling developed extremal length into an effective tool for
geometric function theory:

\begin{itemize}
\tightlist
\item
 They used it to prove the Denjoy--Carleman--Ahlfors theorem.
\item
 They established the \textbf{Ahlfors--Beurling transform} relating
 extremal length to harmonic measure.
\item
 They developed the theory of \textbf{conformal invariants} based on
 extremal length.
\end{itemize}

\subsection{8.4 The Ahlfors--Beurling Conformal
Invariant}\label{ch1936larsvalerianahlfors:the-ahlforsbeurling-conformal-invariant}

The \textbf{extremal distance} between two boundary arcs is defined as
the extremal length of the family of curves connecting them. This gives
a conformal invariant that can be used to estimate harmonic measure,
leading to:

\begin{quote}
\textbf{Ahlfors--Beurling Inequality:} For a simply connected domain
\(\Omega\) and a boundary arc \(E\),
\[\omega(z, E, \Omega) \leq \frac{4}{\pi} \exp(-\pi \lambda(E^c, E))\]
where \(\omega\) is harmonic measure and \(\lambda(E^c, E)\) is the
extremal distance.
\end{quote}

\begin{center}\rule{0.5\linewidth}{0.5pt}\end{center}

\section{9. Conformal Metrics and
Curvature}\label{ch1936larsvalerianahlfors:conformal-metrics-and-curvature}

\subsection{9.1 The Ahlfors Metric}\label{ch1936larsvalerianahlfors:the-ahlfors-metric}

Ahlfors pioneered the use of specially constructed conformal metrics (or
``mass distributions'') to solve problems in complex analysis. A
conformal metric on a domain \(D\) is \(ds = \rho(z)|dz|\) where
\(\rho > 0\) is a smooth function. Its \textbf{Gaussian curvature} is:

\[K_\rho = -\frac{\Delta \log \rho}{\rho^2}\]

\subsection{9.2 The Method of Mass
Distributions}\label{ch1936larsvalerianahlfors:the-method-of-mass-distributions}

In his 1935 paper, Ahlfors introduced the idea of constructing a metric
\(\rho\) on a covering surface such that:

\begin{itemize}
\tightlist
\item
 The metric has curvature bounded above by a constant \(-1\).
\item
 The area of the metric relates to the Nevanlinna characteristic.
\item
 The Gauss--Bonnet theorem gives relations between area, curvature, and
 topology.
\end{itemize}

This method allowed Ahlfors to prove the main theorems of Nevanlinna
theory using only elementary differential geometry.

\subsection{9.3 The Gauss--Bonnet
Version}\label{ch1936larsvalerianahlfors:the-gaussbonnet-version}

In 1937, Ahlfors published a third version of Nevanlinna theory using
the \textbf{Gauss--Bonnet theorem}:

\begin{quote}
\textbf{Ahlfors's Gauss--Bonnet Formula:} For a covering surface
\(\Omega\) with boundary \(\partial\Omega\),
\[\int_\Omega K\, dA + \int_{\partial\Omega} \kappa_g \, ds = 2\pi\chi(\Omega)\]
where \(K\) is the Gaussian curvature and \(\kappa_g\) is the geodesic
curvature of the boundary.
\end{quote}

This version was intended to be generalizable to higher dimensions,
anticipating the later work of Bott, Chern, and Griffiths.

\begin{center}\rule{0.5\linewidth}{0.5pt}\end{center}

\section{10. Kleinian Groups}\label{ch1936larsvalerianahlfors:kleinian-groups}

\subsection{10.1 Definition}\label{ch1936larsvalerianahlfors:definition-2}

A \textbf{Kleinian group} \(\Gamma\) is a discrete subgroup of
\(PSL(2, \mathbb{C})\), acting on the Riemann sphere
\(\overline{\mathbb{C}}\) by Möbius transformations. The limit set
\(\Lambda(\Gamma)\) is the set of accumulation points of orbits; its
complement \(\Omega(\Gamma)\) is the domain of discontinuity, a Riemann
surface (or union of surfaces) on which \(\Gamma\) acts properly
discontinuously.

\subsection{10.2 The Ahlfors Finiteness
Theorem}\label{ch1936larsvalerianahlfors:the-ahlfors-finiteness-theorem}

Ahlfors proved one of the fundamental results in the theory of Kleinian
groups:

\begin{quote}
\textbf{Ahlfors Finiteness Theorem:} If \(\Gamma\) is a finitely
generated Kleinian group, then the quotient \(\Omega(\Gamma)/\Gamma\) is
a finite union of Riemann surfaces of finite type (compact with finitely
many punctures).
\end{quote}

This was a landmark result that established the modern theory of
Kleinian groups. It shows that despite the potential complexity of a
finitely generated Kleinian group, its geometric output is always
finite.

\subsection{10.3 Consequences}\label{ch1936larsvalerianahlfors:consequences}

The finiteness theorem implies: - The domain of discontinuity has
finitely many components. - Each component has a finite fundamental
group. - The group \(\Gamma\) has a finite-sided fundamental polyhedron
in \(\mathbb{H}^3\). - The Ahlfors measure of the limit set is either 0
or the full Lebesgue measure.

\subsection{10.4 Ahlfors's Measure
Conjecture}\label{ch1936larsvalerianahlfors:ahlforss-measure-conjecture}

Ahlfors conjectured that for a finitely generated Kleinian group, the
limit set \(\Lambda(\Gamma)\) either has Lebesgue measure 0 or is the
whole sphere. This was proved by Dennis Sullivan in the 1980s using
ergodic theory.

\begin{center}\rule{0.5\linewidth}{0.5pt}\end{center}

\section{11. The Ahlfors--Bers Theory}\label{ch1936larsvalerianahlfors:the-ahlforsbers-theory}

\subsection{11.1 Teichmüller Theory}\label{ch1936larsvalerianahlfors:teichmuxfcller-theory}

\textbf{Teichmüller theory} studies the space of complex structures on a
Riemann surface up to isotopy. The \textbf{Teichmüller space} \(T(S)\)
of a surface \(S\) is a complex manifold that parametrizes marked
complex structures.

\subsection{11.2 The Bers Embedding}\label{ch1936larsvalerianahlfors:the-bers-embedding}

Ahlfors and Bers, building on Teichmüller's work, developed the modern
theory of Teichmüller spaces. Bers showed that \(T(S)\) can be embedded
as a bounded domain in \(\mathbb{C}^{3g-3}\) (for \(g \geq 2\)) via the
\textbf{Bers embedding} using Schwarzian derivatives.

\subsection{11.3 Ahlfors's
Contributions}\label{ch1936larsvalerianahlfors:ahlforss-contributions-2}

Ahlfors contributed to Teichmüller theory through: - The measurable
Riemann mapping theorem, which shows that every quasiconformal map gives
a point in Teichmüller space. - The study of boundary behavior of
Teichmüller mappings. - The connection between Teichmüller theory and
Kleinian groups.

\begin{center}\rule{0.5\linewidth}{0.5pt}\end{center}

\section{12. Harmonic Analysis on Riemann
Surfaces}\label{ch1936larsvalerianahlfors:harmonic-analysis-on-riemann-surfaces}

\subsection{12.1 The Ahlfors--Sario}\label{ch1936larsvalerianahlfors:the-ahlforssario}

Ahlfors, together with Leo Sario, developed a classification theory of
Riemann surfaces based on the existence of non-constant harmonic
functions. A Riemann surface is classified as: - \textbf{Hyperbolic:} If
it carries a non-constant negative subharmonic function (equivalently, a
Green's function). - \textbf{Parabolic:} If it does not carry a Green's
function (e.g., \(\mathbb{C}\), \(\mathbb{C}^*\)).

\subsection{12.2 The Classification
Problem}\label{ch1936larsvalerianahlfors:the-classification-problem}

Ahlfors proved several criteria for parabolicity: - A Riemann surface is
parabolic if and only if the extremal length of the family of curves
separating a compact set from the boundary is zero. - A covering surface
of a compact Riemann surface is parabolic if and only if its fundamental
group is amenable.

\subsection{12.3 Applications}\label{ch1936larsvalerianahlfors:applications-1}

This classification has applications to: - The type problem for Riemann
surfaces: deciding whether a given surface is hyperbolic or parabolic. -
The theory of harmonic functions on open Riemann surfaces. - The study
of limit sets of Kleinian groups.

\begin{center}\rule{0.5\linewidth}{0.5pt}\end{center}

\section{13. The Schwarzian Derivative}\label{ch1936larsvalerianahlfors:the-schwarzian-derivative}

\subsection{13.1 Definition}\label{ch1936larsvalerianahlfors:definition-3}

The \textbf{Schwarzian derivative} of a locally univalent holomorphic
function \(f\) is:

\[S(f)(z) = \left(\frac{f''(z)}{f'(z)}\right)' - \frac{1}{2}\left(\frac{f''(z)}{f'(z)}\right)^2\]

It satisfies \(S(f \circ g) = (S(f) \circ g) \cdot (g')^2 + S(g)\).

\subsection{13.2 Ahlfors's Work}\label{ch1936larsvalerianahlfors:ahlforss-work}

Ahlfors made important contributions to the theory of Schwarzian
derivatives: - He showed that \(S(f) = 0\) if and only if \(f\) is a
Möbius transformation. - He related the Schwarzian to the curvature of
the metric \(f^*ds_{hyp}\), where \(ds_{hyp}\) is the hyperbolic metric.
- He used the Schwarzian in his study of Teichmüller spaces via the Bers
embedding.

\subsection{13.3 The Schwarzian and
Univalence}\label{ch1936larsvalerianahlfors:the-schwarzian-and-univalence}

Ahlfors proved univalence criteria using the Schwarzian derivative:

\begin{quote}
\textbf{Ahlfors--Weill Theorem:} If
\(|S(f)(z)| \leq \frac{2}{(1-|z|^2)^2}\) for \(|z| < 1\), then \(f\) is
univalent (schlicht) on the unit disc.
\end{quote}

This is the sharp version of a classical theorem of Nehari.

\begin{center}\rule{0.5\linewidth}{0.5pt}\end{center}

\section{14. The Riemann Mapping
Theorem}\label{ch1936larsvalerianahlfors:the-riemann-mapping-theorem}

\subsection{14.1 Classical Statement}\label{ch1936larsvalerianahlfors:classical-statement}

The Riemann mapping theorem states that any simply connected proper
subdomain of \(\mathbb{C}\) is conformally equivalent to the unit disc
\(\mathbb{D}\).

\subsection{14.2 Ahlfors's Proofs}\label{ch1936larsvalerianahlfors:ahlforss-proofs}

Ahlfors gave several proofs of the Riemann mapping theorem, including: -
A proof using extremal length and the \textbf{principle of the
hyperbolic metric}. - A proof using the Ahlfors--Schwarz lemma and
curvature. - A proof using harmonic functions and the Dirichlet
principle.

These proofs exemplify Ahlfors's geometric approach to complex analysis,
each highlighting a different aspect of the subject. The extremal length
proof is particularly elegant: the conformal map is obtained as the
extremal for the problem of maximizing the derivative at a point subject
to boundedness.

\subsection{14.3 The Koebe Distortion
Theorem}\label{ch1936larsvalerianahlfors:the-koebe-distortion-theorem}

A related result that Ahlfors treated extensively is the \textbf{Koebe
distortion theorem}, which bounds the derivative of a univalent
function. Ahlfors gave a proof using the Schwarz lemma and hyperbolic
geometry:

\begin{quote}
\textbf{Koebe Theorem:} If \(f: \mathbb{D} \to \mathbb{C}\) is univalent
with \(f(0) = 0\) and \(f'(0) = 1\), then
\[\frac{1-|z|}{(1+|z|)^3} \leq |f'(z)| \leq \frac{1+|z|}{(1-|z|)^3}\]
\end{quote}

Ahlfors's approach using the hyperbolic metric on the complement of the
image made the proof transparent and unified.

\subsection{14.3 The Boundary
Correspondence}\label{ch1936larsvalerianahlfors:the-boundary-correspondence}

Ahlfors also studied the boundary behavior of conformal maps, proving: -
For a Jordan domain, the conformal map extends to a homeomorphism of the
closure. - The boundary correspondence is Hölder continuous if the
boundary is a quasicircle. - The Carathéodory theorem on prime ends and
its generalization.

\begin{center}\rule{0.5\linewidth}{0.5pt}\end{center}

\section{15. The Ahlfors--Sario Classification of Riemann
Surfaces}\label{ch1936larsvalerianahlfors:the-ahlforssario-classification-of-riemann-surfaces}

\subsection{15.1 Null Classes}\label{ch1936larsvalerianahlfors:null-classes}

Ahlfors and Sario classified Riemann surfaces by their \textbf{null
classes} with respect to various function spaces: - \(O_G\): surfaces
with no non-constant Green's function. - \(O_{HB}\): surfaces with no
non-constant bounded harmonic functions. - \(O_{HP}\): surfaces with no
non-constant positive harmonic functions. - \(O_{HD}\): surfaces with no
non-constant harmonic functions of finite Dirichlet integral.

\subsection{15.2 Relations Between Null
Classes}\label{ch1936larsvalerianahlfors:relations-between-null-classes}

Ahlfors established inclusion relations:
\[O_G \subset O_{HP} \subset O_{HB} \subset \text{(all parabolic surfaces)}\]

and a duality theory between null classes on a surface and its covering
spaces.

\subsection{15.3 The Type Problem}\label{ch1936larsvalerianahlfors:the-type-problem}

The \textbf{type problem} asks whether a given simply connected Riemann
surface is conformally equivalent to \(\mathbb{C}\) (parabolic type) or
\(\mathbb{D}\) (hyperbolic type). Ahlfors proved that the type is
equivalent to the existence of a non-constant bounded harmonic function,
connecting the type problem to potential theory.

\begin{center}\rule{0.5\linewidth}{0.5pt}\end{center}

\section{16. Lattès Maps and Complex
Dynamics}\label{ch1936larsvalerianahlfors:lattuxe8s-maps-and-complex-dynamics}

\subsection{16.1 Lattès Maps}\label{ch1936larsvalerianahlfors:lattuxe8s-maps}

A \textbf{Lattès map} is a rational function
\(R: \overline{\mathbb{C}} \to \overline{\mathbb{C}}\) that is a
quotient of a linear map on a complex torus. Ahlfors studied the
dynamics of such maps and their connection to quasiconformal geometry.

\subsection{16.2 Ahlfors's Contribution}\label{ch1936larsvalerianahlfors:ahlforss-contribution}

Ahlfors showed that Lattès maps provide examples of rational functions
whose Julia set is the whole sphere, and whose dynamics are related to
the measurable Riemann mapping theorem. This work anticipated later
developments in complex dynamics by Douady, Hubbard, and others.

\begin{center}\rule{0.5\linewidth}{0.5pt}\end{center}

\section{17. The Ahlfors Conjecture}\label{ch1936larsvalerianahlfors:the-ahlfors-conjecture}

\subsection{17.1 Statement}\label{ch1936larsvalerianahlfors:statement}

The \textbf{Ahlfors conjecture} (also called the Ahlfors measure
conjecture) states:

\begin{quote}
\textbf{Ahlfors Conjecture:} For a finitely generated Kleinian group
\(\Gamma\), the limit set \(\Lambda(\Gamma)\) either has Lebesgue
measure 0 or is the whole Riemann sphere.
\end{quote}

\subsection{17.2 Proof by Sullivan}\label{ch1936larsvalerianahlfors:proof-by-sullivan}

Dennis Sullivan proved the conjecture in 1982 using ergodic theory on
the geodesic flow on the hyperbolic 3-manifold \(\mathbb{H}^3/\Gamma\).
The key insight: the measure of the limit set is either 0 or 1 by
ergodicity of the geodesic flow, and if it is positive, then the limit
set must be the whole sphere by minimality.

\subsection{17.3 Significance}\label{ch1936larsvalerianahlfors:significance-1}

The Ahlfors conjecture, now a theorem, is fundamental to the modern
theory of Kleinian groups and hyperbolic 3-manifolds. It guarantees that
``small'' limit sets are truly negligible, while ``large'' ones fill the
entire sphere.

\begin{center}\rule{0.5\linewidth}{0.5pt}\end{center}

\section{18. The Ahlfors--Bers Space of Quasiconformal
Mappings}\label{ch1936larsvalerianahlfors:the-ahlforsbers-space-of-quasiconformal-mappings}

\subsection{18.1 The Space of Beltrami
Differentials}\label{ch1936larsvalerianahlfors:the-space-of-beltrami-differentials}

The \textbf{space of Beltrami differentials} \(B(\Gamma)\) on a Riemann
surface \(S\) consists of measurable \((-1,1)\)-forms \(\mu\) with
\(\|\mu\|_\infty < 1\) that are invariant under a Fuchsian group
\(\Gamma\).

\subsection{18.2 The Bers Slice}\label{ch1936larsvalerianahlfors:the-bers-slice}

Bers showed that the space \(B(\Gamma)\) parametrizes the Teichmüller
space \(T(S)\) via the measurable Riemann mapping theorem. The
\textbf{Bers embedding} maps \(T(S)\) into the space of quadratic
differentials, giving it the structure of a bounded domain.

\subsection{18.3 Ahlfors's Role}\label{ch1936larsvalerianahlfors:ahlforss-role}

Ahlfors's work on quasiconformal mappings provided the analytic
foundation for this parametrization. He showed that the dependence of
the quasiconformal map on \(\mu\) is complex-analytic, enabling the use
of analytic function theory in Teichmüller theory.

\begin{center}\rule{0.5\linewidth}{0.5pt}\end{center}

\section{19. The Theory of Harmonic
Integrals}\label{ch1936larsvalerianahlfors:the-theory-of-harmonic-integrals}

\subsection{19.1 Background}\label{ch1936larsvalerianahlfors:background-1}

Ahlfors contributed to the theory of harmonic integrals on Riemann
surfaces, developing the \textbf{Ahlfors--Sario theory} of harmonic
differentials. A harmonic differential \(\omega\) on a Riemann surface
satisfies \(d\omega = 0\) and \(d*\omega = 0\), where \(*\) is the Hodge
star operator.

\subsection{19.2 Ahlfors's Results}\label{ch1936larsvalerianahlfors:ahlforss-results}

Ahlfors proved: - The existence of harmonic differentials with
prescribed periods on open Riemann surfaces. - The \textbf{Ahlfors
decomposition} of a Riemann surface into a compact part and finitely
many ends. - The finiteness of the space of square-integrable harmonic
differentials on a surface of finite type.

\subsection{19.3 Applications}\label{ch1936larsvalerianahlfors:applications-2}

These results are used in: - The theory of moduli of Riemann surfaces. -
The classification of Kleinian groups. - The study of holomorphic
quadratic differentials.

\begin{center}\rule{0.5\linewidth}{0.5pt}\end{center}

\section{20. Meromorphic Functions and Value
Distribution}\label{ch1936larsvalerianahlfors:meromorphic-functions-and-value-distribution}

\subsection{20.1 The Ahlfors Five-Island
Theorem}\label{ch1936larsvalerianahlfors:the-ahlfors-five-island-theorem}

Ahlfors proved a striking generalization of Picard's theorem:

\begin{quote}
\textbf{Ahlfors Five-Island Theorem:} If \(f\) is a meromorphic function
on the plane and \(D_1, \ldots, D_5\) are five simply connected domains
on the Riemann sphere with disjoint closures, then \(f\) maps some
region conformally onto at least one of the \(D_j\).
\end{quote}

This is a purely topological statement about covering surfaces and is
a fundamental result in value distribution theory.

\subsection{20.2 Ahlfors's Island
Theorem}\label{ch1936larsvalerianahlfors:ahlforss-island-theorem}

More generally, Ahlfors proved that if the spherical area of the
complement of \(\bigcup D_j\) is less than \(4\pi\), then there must be
an island. The number 5 is sharp: there exist entire functions that
avoid three given simply connected regions, but not five.

\begin{center}\rule{0.5\linewidth}{0.5pt}\end{center}

\section{21. The Ahlfors--Weill
Extension}\label{ch1936larsvalerianahlfors:the-ahlforsweill-extension}

\subsection{21.1 The Extension Theorem}\label{ch1936larsvalerianahlfors:the-extension-theorem}

Ahlfors and Weill proved a theorem on extending a conformal map of the
unit circle to the interior:

\begin{quote}
\textbf{Ahlfors--Weill Extension Theorem:} If \(\phi\) is a \(C^1\)
function on the unit circle satisfying
\(|\tau \phi''(\tau)| \leq c < 1\) where \(\tau = e^{i\theta}\), then
\(\phi\) extends to a quasiconformal map of the disc.
\end{quote}

\subsection{21.2 Applications}\label{ch1936larsvalerianahlfors:applications-3}

This theorem is fundamental for: - The theory of universal Teichmüller
space. - Boundary value problems for Beltrami equations. - The study of
quasisymmetric maps of the circle.

\begin{center}\rule{0.5\linewidth}{0.5pt}\end{center}

\section{22. The Ahlfors Function}\label{ch1936larsvalerianahlfors:the-ahlfors-function}

\subsection{22.1 Definition}\label{ch1936larsvalerianahlfors:definition-4}

The \textbf{Ahlfors function} of a domain
\(\Omega \subset \overline{\mathbb{C}}\) with \(n\) boundary components
is the extremal function for the problem of maximizing \(|f'(a)|\) among
bounded holomorphic functions \(f: \Omega \to \mathbb{D}\) with
\(f(a) = 0\).

\subsection{22.2 Properties}\label{ch1936larsvalerianahlfors:properties-1}

Ahlfors showed: - The Ahlfors function exists and is unique. - It is a
proper holomorphic map of \(\Omega\) onto \(\mathbb{D}\) of degree equal
to the number of boundary components. - It has a finite number of
critical points.

\subsection{22.3 Construction via Theta
Functions}\label{ch1936larsvalerianahlfors:construction-via-theta-functions}

For a multiply connected domain bounded by \(n\) disjoint circles, the
Ahlfors function can be constructed explicitly using \textbf{theta
functions} and the \textbf{Schottky double} of the domain. The
construction proceeds by considering the Schottky group generated by
reflections in the boundary circles and forming the quotient Riemann
surface.

The Ahlfors function is then expressed as a quotient of theta functions
on the Jacobian variety of the Schottky double. This construction shows
the deep connection between the Ahlfors function and the theory of
abelian functions on Riemann surfaces.

\subsection{22.4 Significance}\label{ch1936larsvalerianahlfors:significance-2}

The Ahlfors function is a fundamental object in geometric function
theory, connecting complex analysis to the topology of planar domains.
It generalizes the Riemann mapping theorem to multiply connected domains
and provides a canonical mapping for domains of arbitrary connectivity.

The function has important applications in: - The theory of conformal
mapping of multiply connected domains. - The computation of harmonic
measures and Green's functions. - The study of moduli spaces of planar
domains. - The theory of holomorphic motions.

\begin{center}\rule{0.5\linewidth}{0.5pt}\end{center}

\section{23. The Ahlfors Laplacian and Harmonic
Analysis}\label{ch1936larsvalerianahlfors:the-ahlfors-laplacian-and-harmonic-analysis}

\subsection{23.1 Ahlfors Regular
Measures}\label{ch1936larsvalerianahlfors:ahlfors-regular-measures}

A \textbf{measure} \(\mu\) on \(\mathbb{R}^n\) is \textbf{Ahlfors
regular} if there exist constants \(C > 0\) and \(d \geq 0\) such that:

\[C^{-1} r^d \leq \mu(B(x, r)) \leq C r^d\]

for all \(x \in \text{supp}(\mu)\) and
\(r \leq \text{diam}(\text{supp}(\mu))\).

\subsection{23.2 Ahlfors's Contribution}\label{ch1936larsvalerianahlfors:ahlforss-contribution-1}

While the formal definition of Ahlfors regularity came later, the
concept is already present in Ahlfors's work on harmonic measure and the
geometry of limit sets. He showed that the harmonic measure on the
boundary of a simply connected domain satisfies such regularity
estimates.

\subsection{23.3 Connection to Kleinian
Groups}\label{ch1936larsvalerianahlfors:connection-to-kleinian-groups}

For a Kleinian group \(\Gamma\) with limit set \(\Lambda\), the
Hausdorff dimension of \(\Lambda\) satisfies a regularity condition that
follows from Ahlfors's finiteness theorem and entropy estimates. The
\textbf{critical exponent} \(\delta(\Gamma) = \dim_H(\Lambda)\) is
bounded by the growth rate of the group and satisfies
\(\delta(\Gamma) \leq 2\) with equality only for groups of the first
kind.

Ahlfors's work on harmonic measure on limit sets established the
fundamental inequality:

\[\dim_H(\Lambda) \leq \inf\{d: \sum_{\gamma \in \Gamma} e^{-d \ell(\gamma)} < \infty\}\]

relating the Hausdorff dimension to the spectral radius of the group.

\subsection{23.4 The Ahlfors Transform}\label{ch1936larsvalerianahlfors:the-ahlfors-transform}

The \textbf{Ahlfors transform} (also called the Beurling--Ahlfors
transform) is the singular integral operator on \(\mathbb{C}\) defined
by:

\[(Tf)(z) = -\frac{1}{\pi} \int_{\mathbb{C}} \frac{f(w)}{(z-w)^2} \, du\,dv\]

where \(w = u + iv\). This operator is the complex form of the Hilbert
transform and satisfies \(\|T\|_p = \max(p, p/(p-1))\) for
\(1 < p < \infty\).

The Ahlfors transform plays a central role in the theory of
quasiconformal mappings: the Beltrami equation
\(\partial_{\bar{z}} f = \mu \partial_z f\) can be solved using the
Ahlfors transform via the singular integral equation
\(f = z + T(\mu f_z)\).

\begin{center}\rule{0.5\linewidth}{0.5pt}\end{center}

\section{\texorpdfstring{24. The Textbook \emph{Complex
Analysis}}{24. The Textbook Complex Analysis}}\label{ch1936larsvalerianahlfors:the-textbook-complex-analysis}

\subsection{24.1 Publication History}\label{ch1936larsvalerianahlfors:publication-history}

Ahlfors's textbook \emph{Complex Analysis} was first published in 1953
by McGraw-Hill, with second and third editions in 1966 and 1979. It has
been translated into numerous languages and remains one of the most
widely used and respected textbooks in the field.

\subsection{24.2 Content and Approach}\label{ch1936larsvalerianahlfors:content-and-approach}

The book covers: - Complex numbers and elementary functions. - Analytic
functions, Cauchy's theorem and integral formula. - Series expansions,
singularities, and the residue theorem. - Conformal mapping and the
Riemann mapping theorem. - Harmonic functions and the Dirichlet problem.
- Elliptic functions and modular functions.

\subsection{24.3 Influence}\label{ch1936larsvalerianahlfors:influence}

The textbook is notable for its elegant geometric approach and its
treatment of the Riemann mapping theorem using the Ahlfors--Schwarz
lemma. It has influenced generations of mathematicians and is often
described as the finest textbook on complex analysis ever written.

The book's structure has been widely imitated: it begins with an
algebraic treatment of complex numbers, proceeds through analytic
function theory, and culminates in geometric function theory and
elliptic functions. Each chapter contains carefully crafted exercises
that extend the theory in non-trivial ways.

\begin{center}\rule{0.5\linewidth}{0.5pt}\end{center}

\section{25. Awards and Recognition}\label{ch1936larsvalerianahlfors:awards-and-recognition}

\textbf{1936:} - Fields Medal (International Mathematical Union) --- one
of the first two awarded, at the ICM in Oslo.

\textbf{1968:} - Wihuri International Prize (Sibelius Prize), Finland
--- awarded for outstanding scientific achievement.

\textbf{1979:} - Wolf Prize in Mathematics --- ``For his outstanding
work in complex analysis.'' The jury's report concluded: ``Everyone
working today in complex analysis is in some sense a student of
Ahlfors.''

\textbf{Elected to:} - Finnish Academy of Sciences. - Royal Swedish
Academy of Sciences. - American Academy of Arts and Sciences. - National
Academy of Sciences (USA). - Norwegian Academy of Science and Letters.

\begin{center}\rule{0.5\linewidth}{0.5pt}\end{center}

\section{26. Selected Major
Publications}\label{ch1936larsvalerianahlfors:selected-major-publications}

\begin{enumerate}
\def\labelenumi{\arabic{enumi}.}
\item
 \textbf{``Zur Theorie der Überlagerungsflächen.''} \emph{Acta
 Mathematica} 65, 157--194, 1935. --- \emph{The theory of covering
 surfaces and the geometric reformulation of Nevanlinna theory, cited
 in the Fields Medal award.}
\item
 \textbf{``Über die asymptotischen Werte der meromorphen Funktionen
 endlicher Ordnung.''} \emph{Acta Societatis Scientiarum Fennicae} 3,
 1930. --- \emph{The Denjoy--Carleman--Ahlfors theorem on asymptotic
 values.}
\item
 \textbf{``Über eine in der neueren Wertverteilungstheorie betrachtete
 Klasse von Riemannschen Flächen.''} \emph{Acta Mathematica} 58,
 375--433, 1932. --- \emph{Geometric Nevanlinna theory.}
\item
 \textbf{``An extension of Schwarz's lemma.''} \emph{Transactions of
 the American Mathematical Society} 43, 359--364, 1938. --- \emph{The
 Ahlfors--Schwarz lemma.}
\item
 \textbf{``Conformal invariants and quasiconformal mappings.''} (with
 A. Beurling) \emph{Acta Mathematica} 83, 101--129, 1950. --- \emph{The
 theory of extremal length.}
\item
 \textbf{``On quasiconformal mappings.''} \emph{Journal d'Analyse
 Mathématique} 3, 1--58, 1954. --- \emph{Foundations of quasiconformal
 mapping theory.}
\item
 \textbf{``Riemann's influence in complex analysis.''} \emph{Wolf Prize
 Lecture}, 1981. --- \emph{A survey of Riemann's legacy.}
\item
 **``Finitely generated Kleinian groups.''\,'' \emph{American Journal
 of Mathematics} 86, 413--429, 1964. --- \emph{The Ahlfors finiteness
 theorem.}
\item
 \textbf{``The structure of finitely generated Kleinian groups.''}
 (with L. Bers) \emph{Acta Mathematica} 122, 1--17, 1969. ---
 \emph{Joint work on Kleinian groups.}
\item
 \textbf{\emph{Complex Analysis.}} McGraw-Hill, 1953 (3rd ed.~1979).
 --- \emph{The classic textbook.}
\end{enumerate}

\begin{center}\rule{0.5\linewidth}{0.5pt}\end{center}

\begin{center}\rule{0.5\linewidth}{0.5pt}\end{center}

\section*{Exercises}
\addcontentsline{toc}{section}{Exercises}

\begin{exercise} Prove that a bounded entire function is constant (Liouville's theorem) and explain how Ahlfors's work extends this to geometric settings. \end{exercise}

\begin{exercise} Compute the hyperbolic metric on the unit disc $\mathbb{D}$ and verify that it is invariant under Möbius transformations. \end{exercise}

\begin{exercise} Show that the Gaussian curvature of the Poincaré metric $ds = \frac{2|dz|}{1-|z|^2}$ on $\mathbb{D}$ is $K = -1$. \end{exercise}

\begin{exercise} Verify the Ahlfors-Schwarz lemma: if $f: \mathbb{D} \to \mathbb{D}$ is holomorphic, then $f^*(ds_{\text{hyp}}) \leq ds_{\text{hyp}}$. \end{exercise}

\begin{center}\rule{0.5\linewidth}{0.5pt}\end{center}

\section*{Summary}
\addcontentsline{toc}{section}{Summary}

Ahlfors revolutionized geometric function theory by applying hyperbolic geometry to complex analysis. His proof of the Riemann mapping theorem using extremal length, and his work on the value distribution theory of meromorphic functions, established the deep connection between complex analysis and differential geometry.

\section{27. References and Further
Reading}\label{ch1936larsvalerianahlfors:references-and-further-reading}

\subsection{Official Sources}\label{ch1936larsvalerianahlfors:official-sources}

\begin{itemize}
\tightlist
\item
 IMU Fields Medal 1936 Announcement:
 https://www.mathunion.org/imu-awards/fields-medal/fields-medals-1936
\item
 Wolf Prize Citation, 1981.
\item
 Harvard University: Lars V. Ahlfors Biography.
\end{itemize}

\subsection{Key Technical References}\label{ch1936larsvalerianahlfors:key-technical-references}

\begin{itemize}
\tightlist
\item
 Ahlfors, L. \emph{Collected Works}. Birkhäuser, 1982. (2 volumes).
\item
 Ahlfors, L. \emph{Complex Analysis}. McGraw-Hill, 3rd ed., 1979.
\item
 Ahlfors, L. \emph{Conformal Invariants: Topics in Geometric Function
 Theory}. McGraw-Hill, 1973.
\item
 Ahlfors, L. and Sario, L. \emph{Riemann Surfaces}. Princeton
 University Press, 1960.
\item
 Ahlfors, L. \emph{Lectures on Quasiconformal Mappings}. Van Nostrand,
 1966.
\end{itemize}

\subsection{Textbooks and Surveys}\label{ch1936larsvalerianahlfors:textbooks-and-surveys}

\begin{itemize}
\tightlist
\item
 Nevanlinna, R. \emph{Analytic Functions}. Springer, 1970.
\item
 Hayman, W. K. \emph{Meromorphic Functions}. Oxford University Press,
 1964.
\item
 Lehto, O. and Virtanen, K. I. \emph{Quasiconformal Mappings in the
 Plane}. Springer, 1973.
\item
 Gardiner, F. P. \emph{Teichmüller Theory and Quadratic Differentials}.
 Wiley, 1987.
\item
 Lehto, O. \emph{Univalent Functions and Teichmüller Spaces}. Springer,
 1987.
\end{itemize}

\subsection{Biography and Exposition}\label{ch1936larsvalerianahlfors:biography-and-exposition}

\begin{itemize}
\tightlist
\item
 Carathéodory, C. ``The Work of Lars Ahlfors.'' ICM 1936 Laudatio.
\item
 Abikoff, W. ``Lars Valerian Ahlfors.'' \emph{Notices of the AMS} 45,
 1998.
\item
 Bers, L. ``Lars Ahlfors: A Biographical Memoir.'' \emph{National
 Academy of Sciences}, 1998.
\item
 Lehto, O. \emph{Lars Ahlfors: Mathematician}. \emph{American
 Mathematical Monthly} 105, 1998.
\item
 Carathéodory, C. ``The First Fields Medals.'' \emph{Proceedings of the
 International Congress of Mathematicians}, Oslo, 1936.
\end{itemize}

\begin{center}\rule{0.5\linewidth}{0.5pt}\end{center}


\begin{center}\rule{0.5\linewidth}{0.5pt}\end{center}

\section*{See Also}
\addcontentsline{toc}{section}{See Also}

This chapter relates to the following work by other Fields Medalists:
\begin{itemize}
  \item \textbf{\hyperref[ch1954kunihikokodaira:kunihiko-kodaira-fields-medal-1954-a-comprehensive-report]{See Chapter on related work}} --- Ahlfors's work on Riemann surfaces relates to Kodaira.
  \item \textbf{\hyperref[ch1958renethom:renuxe9-fruxe9duxe9ric-thom-fields-medal-1958-a-comprehensive-report]{See Chapter on related work}} --- Thom's cobordism relates to Ahlfors.
\end{itemize}

\emph{(End of report)}
